\documentclass[letterpaper]{article}
\usepackage[submission]{aaai2027}
\usepackage[hyphens]{url}
\usepackage{graphicx}
\urlstyle{rm}
\def\UrlFont{\rm}
\usepackage{natbib}
\usepackage{caption}
\usepackage{booktabs}
\usepackage{amsmath}
\usepackage{amssymb}
\usepackage{array}
\frenchspacing

\pdfinfo{
/TemplateVersion (2027.1)
}

\setcounter{secnumdepth}{2}

\title{Supplementary Material for TerraState: A Testable Predictive-State World Model for Weather-Driven Land-Surface Forecasting}
\author{Anonymous Submission}
\affiliations{}

\begin{document}
\maketitle

This technical appendix provides details omitted from Sections~3 and 4 of the
main paper. Appendix~A specifies implementation details of the predictive
state, transition, readout, and training targets; Appendix~B records the
training, preprocessing, and scoring protocol; and Appendices~C--D give the
complete Q2 and Q3 intervention procedures. Numerical Q1--Q3 results remain
in the main paper to avoid duplicating its tables.

\appendix

% Mini-outline: context isolation; state dimensions; transition; training targets.
\section{Additional Implementation Details}
\label{sec:method-details}

\subsection{Context Isolation and Predictive-State Construction}

At forecast issue time \(t\), the optical context is
\(\mathcal H_t=\{(x_i,m_i,\tau_i)\}_{i=1}^{C}\), where \(x_i\) is a
Sentinel-2 composite, \(m_i\) is its validity mask, and \(\tau_i\) is the
acquisition time. Historical meteorology is denoted by
\(u_{\leq t}^{\mathrm{past}}\), the ordered future meteorological sequence by
\(u_{t+1:t+H}\), static geography by \(g\), and the future optical targets by
\(y_{t+1:t+H}\). The context
\(\widetilde{\mathcal C}_t=(\mathcal H_t,u_{\leq t}^{\mathrm{past}},g)\)
contains neither future optical observations nor future meteorology.

Each minicube contains 30 five-day composites at \(128\times128\) pixels. The
first \(C=10\) composites form the context, and the following \(H=20\)
composites form the 100-day prediction window. Optical input has five
channels, including normalized difference vegetation index (NDVI) and four
reflectance bands. Meteorology has 24 channels, obtained from eight variables
with mean, minimum, and maximum aggregation. Static input has five channels;
the first three elevation products provide the patch-wise geography input.

The history operator uses the PVT v2/Contextformer configuration adopted by
TerraState. A \(4\times4\) patch size yields \(32\times32=1{,}024\)
spatial tokens for each minicube. The projector applies layer normalization,
a two-layer multilayer perceptron, and output normalization to the
final-context token at every patch. The resulting per-minicube state has shape
\([1024,256]\), or \([1024B,256]\) after flattening the batch and patch
dimensions.

The context-only constraint is structural. Future optical tensors and future
meteorology are zeroed before the history pass, and future image positions are
masked by the history operator. Consequently, \(b_{1:H}\) and \(z_t\) depend
only on historical observations, past meteorology, and static geography. This
separation makes the later weather substitution specific to the state
transition.

\subsection{Direct Shared Transition and State Readout}

For each queried horizon, the shared transition is conditioned on the
corresponding ordered future-weather prefix. A shared gated recurrent unit
produces codes for all 20 prefixes in one recurrent pass. A geography encoder
average-pools the three elevation channels to the \(4\times4\) patch grid, and
a sinusoidal embedding represents elapsed horizon. The fusion network and
residual transition share parameters across patches and horizons; geography
remains patch-specific, whereas weather and horizon codes are broadcast
across the patch grid.

Every horizon is queried directly from the same history-derived state
\(z_t\). The model applies the shared transition once for each \(h\), using
the prefix \(u_{t+1:t+h}\), rather than feeding \(z_{t+h-1}\) recursively into
the next horizon. The state readout maps each 256-dimensional transitioned token
to one local \(4\times4\) output patch and unpatchifies the patches into
\(r_h\in\mathbb R^{1\times128\times128}\). The coefficient multiplying this
contribution is a fixed non-learnable buffer with value one, which supplies
the precise Q2 cut point.

\subsection{Training-Target Construction}

The ground-truth objective first averages squared error over clear horizons
for each pixel and then averages over eligible vegetation pixels. A frozen
full-weather knowledge-distillation (KD) teacher supplies a stopped forecast
target, weighted by 0.5 in the complete objective.

The future-state objective anchors only the terminal transitioned state. A
training-start frozen copy of the history operator and projector processes the
observed all-frame optical sequence with recorded masks while future weather
is zeroed. The patch mask requires a fully clear terminal \(4\times4\) patch
containing at least one vegetation pixel. Layer-normalized student and target
tokens are compared by cosine distance. Future optical observations create
this stopped training target only; they are absent from the deployed
inference graph. The teacher and target encoder are discarded after training,
and Q2--Q3 introduce no additional objective.

% Mini-outline: optimizer; schedule; batching; hardware; selection.
\section{Training and Evaluation Protocol}
\label{sec:training-evaluation}

\subsection{Training and Implementation Details}

Table~\ref{tab:training} gives the TerraState training configuration. Training
lasts 40 epochs and exactly 14,880 optimizer updates. The implementation uses
AdamW with \(\beta=(0.9,0.999)\), zero weight decay, and gradient clipping at
norm 1.0. The non-history branch uses a base learning rate of
\(3\times10^{-5}\). The history-operator parameter group uses
\(9.9\times10^{-7}\), equal to 0.033 times the branch rate.

The update schedule controls both trainable parameters and
\(\lambda_s\). During the first 20\% of updates, the history operator remains
frozen and \(\lambda_s\) increases linearly from 0 to 0.02. From 20\% through
80\%, the history operator remains frozen and \(\lambda_s=0.02\). During the
final 20\%, only the final Transformer block of the history operator is
unfrozen, and \(\lambda_s=0.01\). The projector, weather encoder, geography
encoder, fusion network, transition, and readout are optimized throughout.

The global batch is 64, realized as eight samples on each of eight GPUs
without gradient accumulation. The shared learning-rate multiplier increases
linearly during the first 300 optimizer updates and then follows cosine decay,
preserving the ratio between the two parameter-group rates.

Model selection uses validation forecasting performance only. Neither Q2,
Q3, nor OOD-t contributes to selection. The selected model is evaluated
without parameter updates on all three questions, thereby keeping the
intervention results outside checkpoint selection and parameter updates.

\begin{table}[t]
\centering
{\small
\setlength{\tabcolsep}{3pt}
\begin{tabular}{@{}ll@{}}
\toprule
Setting & Value\\
\midrule
Epochs / optimizer updates & 40 / 14,880\\
Optimizer & AdamW\\
\(\beta_1,\beta_2\) / weight decay & \(0.9,0.999\) / 0\\
Branch / history learning rate & \(3{\times}10^{-5}\) / \(9.9{\times}10^{-7}\)\\
Learning-rate schedule & 300-update warm-up, cosine decay\\
Global batch & 64 (8 devices \(\times\) 8 samples)\\
Gradient accumulation & none\\
Gradient clipping & global norm 1.0\\
Precision & FP32\\
Hardware & 8 NVIDIA H200 GPUs\\
\bottomrule
\end{tabular}
}
\caption{TerraState training configuration. Learning rates are listed for the
non-history branch and the history-operator parameter group.}
\label{tab:training}
\end{table}

% Mini-outline: splits; preprocessing; masks; scoring outputs.
\subsection{GreenEarthNet Splits and Preprocessing}

GreenEarthNet supplies separate training, validation, and temporal
out-of-distribution partitions. Training optimizes the model, validation
supports checkpoint selection and the validation Q2 analysis, and OOD-t is
held out from selection. The OOD-t evaluation contains 1,904 minicubes. The
validation Q2 scorer contains 952 target minicubes, of which 589 yield paired
metric units after the scorer's eligibility rules.

For each minicube, the data adapter samples the optical record every five days
to obtain 30 composites. NDVI is computed from near-infrared and red
reflectance. Missing optical values are set to zero after a separate cloud and
quality mask is constructed. The cloud mask combines the provided learned
mask with valid scene-class labels. Meteorological variables are normalized
using the GreenEarthNet statistics and aggregated over each five-day interval
by mean, minimum, and maximum, producing 24 channels. Elevation channels are
scaled by 500, and categorical static layers retain their original scale.

The forecast target is NDVI over valid vegetation pixels. Clear target
observations satisfy the official cloud and scene-class rule, and vegetation
pixels fall in the declared land-cover range. Prediction tensors have shape
\([B,20,1,128,128]\). For file-based scoring, each output stores the 20 NDVI
predictions on the target timestamps with the original latitude and longitude
coordinates. The official scorer computes pixel-level time-series quantities
before balancing them across minicubes and vegetation land-cover classes. It
reports \(R^2\), RMSE, NSE, absolute bias, and \(\mathrm{RMSE}_{25}\), as
defined under \emph{Experimental Setup} in Section~4 of the main paper.

% Mini-outline: arms; global estimand; paired estimand; uncertainty; invariants.
\section{Q2 Intervention and Statistical Protocol}
\label{sec:q2-results}

Q2 evaluates three forward computations on identical samples. The full model
uses the additive state path in the main paper. State removal sets the state
contribution to zero immediately before addition, giving
\(\widehat y^{\mathrm{remove}}_{t+h}=b_h\). The identity-transition control
replaces the learned transition output by the unevolved \(z_t\) while keeping
the state readout active. Model weights, history, target window, masks, and
scorer remain fixed for every arm.

The analysis reports two distinct estimands. Official
\(\Delta R^2\) is the difference between full and intervened dataset-level
scores. Paired \(\Delta R^2\) is computed per eligible minicube and then
averaged; its 95\% percentile interval comes from 10,000 paired bootstrap
resamples. Because the two quantities use different aggregation orders, the
paired interval characterizes the mean paired-minicube effect rather than the
official dataset-level difference.

The state-removal comparison provides the primary evidence: the explicit
state contribution yields a measurable forecast increment on both validation
and OOD-t. The identity-transition arm provides complementary evidence for
the learned state update and is treated as supporting rather than primary
evidence. Numerical effects for both splits are reported in Table~2 of the
main paper.

Implementation invariants check that all arms use identical sample order,
target tensors, masks, history-derived context forecast, and frozen
model weights. They also verify that state removal recovers \(b_h\), the fixed
state-contribution coefficient is restored after each call, and finite paired
units are aligned before bootstrapping.

% Mini-outline: subset; controls; fixed quantities; loss; bootstrap; scope.
\section{Q3 Heat--Drought Subset and Weather-Control Protocol}
\label{sec:q3-results}

Q3 uses a frozen heat--drought subset with 84 matched pairs. The 84 evaluation
samples comprise all eligible OOD-t minicubes in the predeclared
heat--drought stratum. Hot and dry anomaly thresholds were estimated from the
training climatology at the 80th percentile and frozen before model
evaluation; eligible samples additionally required at least 80\% valid
future-weather coverage. Neither future NDVI, model predictions, forecast
errors, nor checkpoint outputs entered subset construction. Donors were
selected from the normal-weather pool by deterministic nearest-neighbor
matching within the same meteorological season. The distance was the
root-mean-square standardized difference over day of year, latitude,
longitude, vegetation fraction, context cloud ratio, context valid fraction,
context NDVI mean, and future-weather valid fraction. Matches were required
to fall within a 1.5 caliper, and each donor could be reused at most five
times. This frozen procedure produced 84 pairs involving 45 unique donors and
31 geographic clusters.

For the weather intervention, each evaluated minicube retains its historical
context, \(b_{1:H}\), \(z_t\), geography, horizon queries, readout, target
window, and mask. Only the future-weather tensor supplied to \(T_\psi\)
changes among three arms:
\begin{enumerate}
    \item actual future weather from the evaluated minicube;
    \item donor future weather from the frozen matching procedure;
    \item normalized-mean weather, represented by zero in the frozen
    global standardized feature space.
\end{enumerate}

The primary loss is masked mean squared error over the complete 20-step
forecast window. Control-minus-actual loss is computed per pair, so a positive
value favors actual weather. The uncertainty analysis resamples the 31
geographic clusters with replacement for 10,000 replicates, preserving
dependence caused by donor reuse. Whether actual weather yields lower loss is
summarized by a descriptive pair count, while statistical inference relies on
the geographic-cluster bootstrap. The subset \(R^2\) and RMSE in Table~3 of
the main paper refer specifically to these 84 matched samples.

This intervention measures conditional response fidelity under a frozen
matched protocol. This estimand characterizes response fidelity under observed
targets; causal and counterfactual identification lie outside its scope
because alternate real-world outcomes under the control weather sequences are
unobserved.

\end{document}
